% Generated draft: requires mathematical and source review.
\documentclass[11pt]{article}
\usepackage[margin=25mm]{geometry}
\usepackage{fontspec}
\setmainfont{Latin Modern Roman}
\usepackage{amsmath,amssymb,mathtools,unicode-math}
\setmathfont{Latin Modern Math}
\usepackage{enumitem,needspace,xurl,hyperref}
\setlength{\parindent}{0pt}
\setlength{\parskip}{5pt}
\setlength{\emergencystretch}{3em}
\newcommand{\N}{\mathbb{N}}
\newcommand{\Z}{\mathbb{Z}}
\newcommand{\Q}{\mathbb{Q}}
\newcommand{\R}{\mathbb{R}}
\newcommand{\C}{\mathbb{C}}
\newcommand{\F}{\mathbb{F}}
\newcommand{\A}{\mathbb{A}}
\newcommand{\PP}{\mathbb{P}}
\newcommand{\E}{\mathbb{E}}
\newcommand{\eps}{\varepsilon}
\newcommand{\dd}{\,\mathrm{d}}
\begin{document}
\section*{UnsolvedMath ID: 1875}
\textbf{Source problem code:} catalogue entry
\subsection{Definitions and mathematical statement}
A base-2 pseudoprime is an odd composite integer \(N>1\) such that
\[
2^{N-1}\equiv1\pmod N.
\]
A perfect square has the form \(N=a^2\) with \(a\in\Z\), \(a\geq2\). A non-squarefree integer instead only requires some prime \(q\) with \(q^2\mid N\); it need not itself be a square. Set
\[
Q_1=1093^2=1194649,\qquad Q_2=3511^2=12327121.
\]
The literal perfect-square reading of the supplied question is the existence problem
\[
\exists a\geq2:\quad N=a^2\text{ is a base-2 pseudoprime},\quad
Q_1\nmid N,\quad Q_2\nmid N.
\]
Both divisibility exclusions are retained: ``other than multiples'' is not weakened to merely \(N\ne Q_1,Q_2\).

The background also discusses all non-squarefree pseudoprimes. That alternative replaces the condition \(N=a^2\) by
\[
\exists q\text{ prime}\quad q^2\mid N,
\]
while retaining the pseudoprime condition and the two divisibility exclusions. This is a different, broader existence problem. The source does not settle which scope was intended.

An odd prime \(p\) is a base-2 Wieferich prime if \(2^{p-1}\equiv1\pmod{p^2}\). The supplied triage dated 2026-08-17 records the equivalence
\[
p^2\text{ is a base-2 pseudoprime}
\quad\Longleftrightarrow\quad
2^{p-1}\equiv1\pmod{p^2}.
\]
The prime-square subcase of the displayed existence problem therefore asks for a Wieferich prime other than 1093 and 3511. The triage cites Dorais--Klyve's certified search showing no others for \(p<6.7\times10^{15}\); this is a bound on the prime \(p\), not on every square \(N\), and does not classify squares of composite bases or all non-squarefree pseudoprimes globally.

The background attributes to Pinch a list of 54 non-squarefree base-2 pseudoprimes up to \(10^{13}\), all divisible by \(Q_1\) or \(Q_2\), but supplies no separate citation for that exact computation. The registry labels the entry open. The dated triage labels it partially solved and reports no complete global classification for either natural scope. These are supplied historical and dated status statements, not new verification.
\subsection{Short English statement}
Is there a perfect-square base-2 pseudoprime divisible by neither \(1093^2\) nor \(3511^2\)? The source also discusses the broader non-squarefree question. The known certified search concerns prime squares via Wieferich primes and does not settle either global classification.
\subsection{Sources}
\begin{itemize}
\item Local UnsolvedMath registry, permanent ID 1875, source code catalogue entry, ``Square Pseudoprimes,'' including supplied literature triage dated 2026-08-17. The background attributes the question to Richard Guy, \emph{Unsolved Problems in Number Theory}, Section A12, and the historical list to Pinch; no edition, page, or separate Pinch reference is supplied.
\item Fran\c{c}ois G. Dorais and Dominic Klyve, ``A Wieferich prime search up to \(6.7\times10^{15}\),'' \emph{Journal of Integer Sequences} 14 (2011), Article 11.9.2: \url{https://cs.uwaterloo.ca/journals/JIS/VOL14/Klyve/klyve3.html}.
\item Carl Pomerance, J. L. Selfridge, and Samuel S. Wagstaff Jr., ``The pseudoprimes to \(25\times10^9\),'' \emph{Mathematics of Computation} 35 (1980), 1003--1026: \url{https://doi.org/10.1090/S0025-5718-1980-0572872-7}.
\end{itemize}
\end{document}
